\documentclass[12pt]{article}
\usepackage{palestra1}
\begin{document}
\begin{ficha}{Sumar y restar}{Entrar}{2}

\begin{encargo}
En la palestra, el \textbf{+} quiere decir «que entre». La bandera acaba en la
suma de los que tiran.
\end{encargo}

\bloque{1}{Mira cómo se hace}{En el campo}\quad{\small así se juega $-3 + 5$}

\vspace{1.5mm}
\begin{minipage}[c]{.5\linewidth}
\paso{1}\textbf{Entra el \nm{-3}.} La bandera, en \nm{-3}.\\[1.5mm]
\paso{2}\textbf{Entra el \nm{5}.} Tira más fuerte.\\[1.5mm]
\paso{3}\textbf{¿Dónde acaba?} En \nm{2}:\ \ $-3 + 5 = +2$
\end{minipage}%
\begin{minipage}[c]{.5\linewidth}\centering
\campo[paso=3.2mm,escala=.6]{-3,5}
\end{minipage}

\vspace{1.5mm}
\begin{listillon}
Si los dos tienen el \textbf{mismo signo}, se suman las fuerzas: $-2 + (-3) = -5$.
Si tienen \textbf{distinto signo}, se restan, y gana el signo del que tira más
fuerte: $-4 + 6 = +2$.
\end{listillon}

\bloque{2}{Ahora tú}{En el campo}\quad{\small escribe la cuenta y dónde acaba la bandera}

\vspace{1.5mm}
\begin{tabular}{@{}c@{\hspace{8mm}}c@{}}
\campo[bandera=?,paso=3.2mm,escala=.6]{-4,6} & \campo[bandera=?,paso=3.2mm,escala=.6]{-2,-3}\\[1mm]
\ej{a}\hueco[3cm] $=$ \casillita[8mm] & \ej{b}\hueco[3cm] $=$ \casillita[8mm]\\
\end{tabular}

\bloque{3}{Sin campo}{Sobre el papel}

\vspace{1mm}
{\renewcommand{\arraystretch}{2.1}
\begin{tabular}{@{}l@{\hspace{14mm}}l@{}}
\ej{a}$-5 + 2 = $ \casillita[8mm] & \ej{b}$3 + (-7) = $ \casillita[8mm]\\
\ej{c}$-1 + (-4) = $ \casillita[8mm] & \ej{d}$6 + (-6) = $ \casillita[8mm]\\
\end{tabular}}

\alto{$-2 + (-3)$ no es $+5$: los dos tiran hacia la izquierda, y la bandera se
va a $-5$.}

\reto{Haz que la bandera acabe en el 0 con dos tiradores. ¿De cuántas maneras se
puede?}

\comprueba{Cada código abre esa cuenta en Practicar.}{%
\qr{1}{j=sumar\&m=2\&c=n-3_s_n5}\hfill\qr{2a}{j=sumar\&m=2\&c=n-4_s_n6}\hfill
\qr{2b}{j=sumar\&m=2\&c=n-2_s_n-3}\hfill\qr{3a}{j=sumar\&m=2\&c=n-5_s_n2}\hfill
\qr{3b}{j=sumar\&m=2\&c=n3_s_n-7}\hfill\qr{3c}{j=sumar\&m=2\&c=n-1_s_n-4}}

\end{ficha}

% ─────────────────────────────── REVERSO ──────────────────────────────────
\begin{dorso}{Sumar y restar}{Retirarse}{2}

\bloque{1}{Mira cómo se hace}{En el campo}\quad{\small el \textbf{−} es «que se retire»: así se juega $2 - (-3)$}

\vspace{1mm}
{\small
\begin{tabular}{@{}c@{\hspace{3mm}}c@{\hspace{3mm}}c@{}}
\campo[alcance=6,paso=3.2mm,escala=.5]{2} &
\campo[alcance=6,paso=3.2mm,escala=.5]{2,-3,3} &
\campo[alcance=6,paso=3.2mm,escala=.5]{2,3}\\
\parbox[t]{5.2cm}{\centering\paso{1}Tira el \nm{2}.} &
\parbox[t]{5.2cm}{\centering\paso{2}No hay ningún \nm{-3}: entra con su pareja, el \nm{3}. Suman 0.} &
\parbox[t]{5.2cm}{\centering\paso{3}Se retira el \nm{-3}. La bandera, en \nm{5}.}\\
\end{tabular}}

\vspace{1.5mm}
\begin{listillon}
\textbf{Restar} es \textbf{sumar el opuesto}: $2 - (-3) = 2 + 3 = 5$.
\end{listillon}

\bloque{2}{Ahora tú}{En el campo}

\vspace{1mm}
{\renewcommand{\arraystretch}{1.9}
\begin{tabular}{@{}l@{\hspace{8mm}}c@{\hspace{8mm}}c@{\hspace{8mm}}c@{}}
& \small\textbf{¿está?} & \small\textbf{¿quién entra con su pareja?} & \small\textbf{bandera}\\
\ej{a}$4 - (+1)$ & \hueco[1.2cm] & \hueco[2.6cm] & \casillita[8mm]\\
\ej{b}$1 - (-2)$ & \hueco[1.2cm] & \hueco[2.6cm] & \casillita[8mm]\\
\ej{c}$-3 - 2$ & \hueco[1.2cm] & \hueco[2.6cm] & \casillita[8mm]\\
\end{tabular}}

\bloque{3}{Sin campo}{Sobre el papel}

\vspace{1mm}
{\renewcommand{\arraystretch}{2.1}
\begin{tabular}{@{}l@{\hspace{14mm}}l@{}}
\ej{a}$5 - (-1) = $ \casillita[8mm] & \ej{b}$-2 - 4 = $ \casillita[8mm]\\
\end{tabular}}

\vspace{2mm}
\begin{semaforo}
\sfila{Sumo dos números con el mismo signo}
\sfila{Sumo dos números con distinto signo}
\sfila{Resto: sé quién entra con su pareja}
\end{semaforo}

\comprueba{Cada código abre esa cuenta en Practicar.}{%
\qr{1}{j=sumar\&m=3\&c=n2_r_n-3}\hfill\qr{2a}{j=sumar\&m=3\&c=n4_r_n1}\hfill
\qr{2b}{j=sumar\&m=3\&c=n1_r_n-2}\hfill\qr{2c}{j=sumar\&m=3\&c=n-3_r_n2}\hfill
\qr{3a}{j=sumar\&m=3\&c=n5_r_n-1}\hfill\qr{3b}{j=sumar\&m=3\&c=n-2_r_n4}}
\end{dorso}
\end{document}
